Solar irradiance on segmented roof surfaces was estimated using voxel-based ray marching
combined with the Beer--Lambert Law. Roof surfaces were first extracted from the Actueel
Hoogtebestand Nederland third edition (AHN3) airborne Light Detection and Ranging (LiDAR)
point cloud via planar segmentation, with each resulting surface patch assigned a set of
sample points for irradiance evaluation. The surrounding scene --- including buildings and
vegetation --- was rasterised into a 0.5~m isotropic voxel grid covering the full extent
of the study area. For each sample point, rays were cast toward 145 discrete directions
distributed uniformly across the upper hemisphere, following a parameterisation of the sky
dome based on Tregenza subdivision \citep{Tregenza1983}. Along each ray, the cumulative
transmittance $\tau$ was computed by integrating the extinction coefficient $\kappa$ over
the ray path through the voxel grid,

\begin{equation}
    \tau = \exp\!\left(-\int_0^{L} \kappa\!\left(\mathbf{x}(s)\right)\mathrm{d}s\right),
    \label{eq:beer-lambert}
\end{equation}

\noindent where $\mathbf{x}(s)$ is the ray position at path length $s$ and $L$ is the
total ray length to the sky dome boundary. Building voxels were treated as fully opaque
($\kappa \to \infty$, implemented as hard blocking), while vegetation voxels were assigned
an attenuation coefficient of $\kappa = 0.5~\mathrm{m}^{-1}$, consistent with empirical
estimates of plant-area-density extinction for urban tree canopies
\citep{Jonckheere2004,Widlowski2006}. The per-direction transmittance values were then
weighted by the corresponding sky radiance distribution --- using the all-sky Perez model
for diffuse irradiance and direct-beam geometry for beam irradiance --- and integrated to
yield the plane-of-array (POA) irradiance at each sample point. Occlusion by surrounding
structures is therefore fully accounted for through the ray transmittance rather than
through a binary sky view factor, enabling a continuous representation of partial shading
by vegetation.
